% Generated draft: requires mathematical and source review.
\documentclass[11pt]{article}
\usepackage[margin=25mm]{geometry}
\usepackage{fontspec}
\setmainfont{Latin Modern Roman}
\usepackage{amsmath,amssymb,mathtools,unicode-math}
\setmathfont{Latin Modern Math}
\usepackage{enumitem,needspace,xurl,hyperref}
\setlength{\parindent}{0pt}
\setlength{\parskip}{5pt}
\setlength{\emergencystretch}{3em}
\newcommand{\N}{\mathbb{N}}
\newcommand{\Z}{\mathbb{Z}}
\newcommand{\Q}{\mathbb{Q}}
\newcommand{\R}{\mathbb{R}}
\newcommand{\C}{\mathbb{C}}
\newcommand{\F}{\mathbb{F}}
\newcommand{\A}{\mathbb{A}}
\newcommand{\PP}{\mathbb{P}}
\newcommand{\E}{\mathbb{E}}
\newcommand{\eps}{\varepsilon}
\newcommand{\dd}{\,\mathrm{d}}
\begin{document}
\section*{UnsolvedMath ID: 10400162}
\textbf{Source problem code:} catalogue entry
\subsection{Definitions and mathematical statement}
T. Kerler's programme starts with the source's algebraic category \(\operatorname{Alg}\), distinguished Hopf-algebra object \(A\), and functor \(F:\operatorname{Alg}\to\operatorname{Cob}_\bullet\) to its three-dimensional cobordism category. These categories and the functor are not defined in the excerpt. The four tasks are:
\begin{enumerate}
\item Find additional relations on \(\operatorname{Alg}\), beyond the Hopf-algebra axioms, making the induced functor an isomorphism \(\operatorname{Alg}/R\to\operatorname{Cob}_\bullet\). The source asks for a full category isomorphism, not merely a nonfaithful representation.
\item Find relations making the maps
\[
 \operatorname{Aut}_{\operatorname{Alg}/R}(A^{\otimes g})
 \longrightarrow\Gamma_g\cong
 \operatorname{Aut}_{\operatorname{Cob}_\bullet}(\Sigma_{g,1})
\]
isomorphisms. Here \(\Sigma_{g,1}\) is the genus-\(g\) surface with one boundary component, and \(\Gamma_g\) is the source's mapping class group, conventionally with boundary fixed. This automorphism-only problem is presented as easier than recovering all cobordism morphisms.
\item Relate these presentation issues to obstruction groups, such as Steinberg and Whitehead groups, through stratified function spaces.
\item Identify analogous algebraic structures in higher dimensions. The source anticipates higher categories and more generators than one object \(A\), with partial answers potentially leading to functorial four-manifold invariants by linear representations; no particular higher-category model or invariant is specified.
\end{enumerate}
The geometric motivation is that algebraic generators represent Morse-elementary cobordisms, while relations correspond to handle slides and cancellations. The suggested comparison with geometrization and the Poincar\'e conjecture is a dated 2002 analogy, not a present-status claim.

The appended context reports the following result of source reference [216]. An ETQFT, extended to surfaces with boundaries, is a double functor from a double category \(\operatorname{Cob}_*\) of relative, 2-framed, \(1+1+1\)-dimensional cobordisms to a double category of linear abelian categories over a perfect field. Its circle category is \(\mathcal C_V=V(S^1)\), viewing the circle as a 0-object. The source reports a construction \(V_{\mathcal C}\) with \(\mathcal C_{V_{\mathcal C}}=\mathcal C\) for its modular tensor categories, described as bounded ribbon braided tensor categories with additional properties. It reports the semisimple construction and nonsemisimple extensions for connected surfaces with boundary and disconnected closed surfaces, using half-projective TQFTs. These qualifications are retained; \(\operatorname{Cob}_*\) is not identified silently with \(\operatorname{Cob}_\bullet\), and this contextual result does not answer all four tasks.
\subsection{Short English statement}
Find relations presenting three-dimensional cobordisms and their mapping class groups through Kerler's Hopf-algebra category, relate them to obstruction groups, and seek higher-dimensional analogues. The source's extended-TQFT context and its semisimple/nonsemisimple scope remain explicit.
\subsection{Sources}
\begin{itemize}
\item Local UnsolvedMath registry, permanent ID 10400162, source code catalogue entry; statement attribution T. Kerler; proposer field and proposal year absent; update 2026-08-06.
\item Ohtsuki, \emph{Problems on invariants of knots and 3-manifolds} (2002), Problem 8.16, PDF page 131: \url{https://msp.org/gtm/2002/04/gtm-2002-04-024s.pdf}. Registry access date 2026-07-29; PDF font-structure extraction. Provenance distinguishes 231 numbered bold headings. Present primary-literature status and redistribution terms require review.
\item The appended ETQFT construction is credited to source reference [216], whose bibliographic entry is absent. The research summary only identifies Kerler's programme and provides no exact resolving reference.
\end{itemize}
\end{document}
